\section*{Chapitre \ref{ch:m2:repo-and-specials} --- Repo et specials}

\begin{solution}{exo:m2:repo-and-specials:1}
Du vendredi au mardi, il y a quatre jours~: $50\,000\,000 \times 0.039 \times 4/360 =
\text{USD}~21\,666.67$.
\end{solution}

\begin{solution}{exo:m2:repo-and-specials:2}
Valeur du collatéral $100\,000\,000 \times 1.015 = 101\,500\,000$~; moins 2~\%~:
99\,470\,000~USD.
\end{solution}

\begin{solution}{exo:m2:repo-and-specials:3}
\omterm{def:m2:repo-and-specials:special}{Specialness} $3.90 - 3.10 = 80$ points de base. Financer 100 millions de dollars à
3.10~\% au lieu de 3.90~\% permet d'économiser $100\,000\,000 \times 0.008/360 =
\text{USD}~2\,222.22$ par jour (sur les liquidités empruntées~; il faut multiplier par le prix plein coupon pour
obtenir le montant exact).
\end{solution}

\begin{solution}{exo:m2:repo-and-specials:4}
Coupon $100\,000\,000 \times 0.02125 \times 30/184 = 346\,467$~; repo $100\,870\,911
\times 0.039 \times 30/360 = 327\,830$~; carry de 18\,637~USD. DV01 de 80\,444~USD par
point de base~: point mort de 0.23 points de base. À 4.30~\%, le repo coûte 361\,454 et
le carry vaut $-\text{USD}~14\,987$~: le titre rapporte désormais moins que son financement,
et le détenir revient à parier sur une baisse des rendements.
\end{solution}

\begin{solution}{exo:m2:repo-and-specials:5}
À 1~\%~: $\frac{1}{360} \times 0.01 \times 2 \times 100\,000\,000 = 5\,555.56$ par
jour, soit 27\,777.78~USD sur cinq jours. À 3.75~\%~: $\max(3 - 3.75, 0) = 0$~; rien.
\end{solution}

\begin{solution}{exo:m2:repo-and-specials:6}
Avec $R = 0$~: floor $-(3\% - 0) = -3\%$~; les vendeurs à découvert paieraient jusqu'à 3~\% par an pour
emprunter une émission rare avant de laisser la livraison en suspens. Aujourd'hui, avec $R = 3.75\%$~: floor 0.
Avant 2009, il n'existait aucune \omterm{def:m2:repo-and-specials:fails}{pénalité de suspens}~: le défaut de livraison ne coûtait que l'intérêt manqué sur
le produit de la vente, si bien qu'aucun vendeur à découvert ne prêtait de liquidités à un taux inférieur à 0~\%, et lorsque les taux se sont effondrés
à la fin de 2008, ce coût s'est lui aussi évanoui~; les suspens sont devenus l'alternative la moins chère
à l'emprunt, et ils se sont multipliés.
\end{solution}

\begin{solution}{exo:m2:repo-and-specials:7}
Après une baisse de 1.5 points le premier jour~: le prêteur doit recevoir $98\,853\,492$
plus un jour d'intérêts (10\,709), ce qui exige un collatéral valant
$98\,864\,202/0.98$~; le collatéral vaut 99\,370\,911~: un appel de marge de
1\,510\,928~USD. Après 30 jours à prix inchangé~: le seul coupon couru,
majoré de la décote, $321\,274/0.98$~: 327\,830~USD. Dans la pratique, les repos
règlent en général les intérêts à l'échéance et la marge sur le principal~; la
différence tient aux conditions de chaque contrat.
\end{solution}

\begin{solution}{exo:m2:repo-and-specials:8}
Les décisions du Comité n'ont pas changé le 17 septembre~: la fourchette est restée à
2.00--2.25~\% et les taux administrés avec elle. Le taux de marché est sorti de la
fourchette, car les réserves étaient devenues rares pendant une journée (paiements d'impôts et
le règlement des coupons a retiré 120 milliards), de sorte que les banques à court de liquidités ont payé plus cher et celles qui en avaient
les ont prêtées là où les taux étaient les plus élevés, en repo. La Fed a réagi en
\emph{ajoutant} des réserves au moyen d'opérations de repo, et le lendemain elle a abaissé la
fourchette d'un quart de point~: l'inverse d'un resserrement.
\end{solution}

\begin{solution}{pb:m2:repo-and-specials:1}
\textbf{1.} 3.45~\%, 3.65~\% et 3.80~\%.
\textbf{2.} $100\,000\,000 \times 1.00870911 \times 0.0045 \times 30/360 =
\text{USD}~37\,827$.
\textbf{3.} $100\,870\,911 \times (30 \times 0.0045 + 30 \times 0.0025 + 30 \times
0.0010)/360 = \text{USD}~67\,247$.
\textbf{4.} La nouvelle émission est l'instrument de couverture du marché, si bien que les positions courtes sont
les plus importantes juste après l'émission, alors que l'offre de titres que les détenteurs sont prêts à prêter
est encore en train de se constituer.
\textbf{5.} Ce sont les vendeurs à découvert, à travers les liquidités qu'ils prêtent au taux bas pour emprunter le
titre. Ils paient car un défaut de livraison leur coûterait l'intérêt sur le produit de leur vente
(et, lorsque les taux sont bas, la \omterm{def:m2:repo-and-specials:fails}{pénalité de suspens}), et parce que le titre est
la couverture dont ils ont besoin.
\textbf{6.} $V = 100.870911 \times 0.24/360 = 0.0672$ points pour 100.
\textbf{7.} $0.0672 \times 32 = 2.15$ 32es.
\textbf{8.} $0.0672/0.080444 = 0.84$ points de base.
\textbf{9.} Environ $4.210 - 0.008 = 4.202\%$.
\textbf{10.} À 4.195~\%, il est plus cher de 0.7 points de base que ne le justifie le
financement seul~: le reste est le prix de la liquidité.
\textbf{11.} Acheter l'off-the-run et vendre l'on-the-run pour la même DV01~; mettre en repo
l'off-the-run au collatéral général pour le financer~; prendre en \omterm{def:m2:repo-and-specials:rate}{reverse repo}
l'on-the-run (prêter des liquidités, emprunter le titre) pour livrer la vente à découvert~; conserver jusqu'à ce que la
prochaine adjudication rende l'on-the-run ancien.
\textbf{12.} La position longue paie le taux du collatéral général (GC), soit 3.90~\%, sur son financement~; la position courte ne touche que
le \omterm{def:m2:repo-and-specials:special}{taux de repo special} sur les liquidités qu'elle prête, 3.45~\% puis 3.65~\% puis 3.80~\%~: cette
opération paie la \omterm{def:m2:repo-and-specials:special}{specialness}, qui est précisément ce qu'elle parie être cotée trop cher.
\textbf{13.} La position courte paie deux fois plus~: 0.1345 points au lieu de 0.0672, et
l'opération perd, sauf si le spread de rendement joue davantage en sa faveur.
\textbf{14.} Un \omterm{def:m2:repo-and-specials:special}{taux de repo special} de 0~\% porte la \omterm{def:m2:repo-and-specials:special}{specialness} à 3.90~\%~: la position courte paie
10\,928~USD par jour d'intérêts manqués sur 100 millions de dollars, et davantage si un short squeeze
rend le titre impossible à emprunter du tout.
\textbf{15.} Zéro (\cref{prop:m2:repo-and-specials:floor})~: avec le taux de référence
supérieur à 3~\%, la \omterm{def:m2:repo-and-specials:fails}{pénalité de suspens} est nulle et le défaut de livraison ne coûte pas plus cher
qu'un prêt à 0~\%.
\textbf{16.} Liquidité~: l'indice de référence se négocie aux fourchettes les plus étroites, est
accepté partout comme collatéral et comme instrument de couverture, et les investisseurs paient pour cela
au-delà de sa valeur de financement.
\textbf{17.} Si un emprunteur ne le restitue pas, le prêteur perd sa propre capacité
à livrer~; la \omterm{def:m2:repo-and-specials:fails}{pénalité de suspens} est ce que paierait la partie défaillante, et aux
taux actuels elle est nulle, si bien que la protection a disparu.
\textbf{18.} Davantage de titres du même type à prêter~: la \omterm{def:m2:repo-and-specials:special}{specialness} baisse.
\textbf{19.} Résultat à retenir~: \emph{la valeur du special} sur le cycle
d'adjudication est de 0.0672 points pour 100 (2.15 32es), 0.84 points de base de rendement, 67\,247~USD
sur 100 millions de dollars.
\textbf{20.} Un titre tellement recherché à l'emprunt que les prêteurs de liquidités
acceptent un taux inférieur au taux général pour l'obtenir.
\end{solution}

\begin{solution}{iq:m2:repo-and-specials:1}
Une vente de titres assortie d'un engagement de rachat~: un prêt de liquidités contre
collatéral, ou de collatéral contre liquidités. Toute position longue sur obligation est
financée par elle, et toute position courte y est empruntée~; son taux est le
coût de portage des traders, et sa \omterm{def:m2:repo-and-specials:special}{specialness} est le prix des titres dont ils sont
vendeurs à découvert.
\iqlookfor{les deux sens de l'opération, et le financement comme lien vers chaque
position.}
\end{solution}

\begin{solution}{iq:m2:repo-and-specials:2}
Son \omterm{def:m2:repo-and-specials:special}{taux de repo special} est inférieur à celui du collatéral général~: les prêteurs de liquidités acceptent un
taux plus bas pour obtenir ce titre, généralement parce qu'il est demandé pour couvrir
des positions courtes ou pour la livraison. Le détenteur y gagne~: il emprunte des liquidités à meilleur compte
en le donnant en garantie, ou il prête le titre~; ce sont les vendeurs à découvert qui paient.
\iqlookfor{la \omterm{def:m2:repo-and-specials:special}{specialness} comme transfert des vendeurs à découvert vers les détenteurs.}
\end{solution}

\begin{solution}{iq:m2:repo-and-specials:3}
Le 16 septembre 2019, les paiements d'impôts et le règlement-livraison d'un coupon ont retiré environ 120
milliards de réserves, qui sont tombées à 1.34~billion de dollars~; le 17, le \omterm{def:m2:repo-and-specials:rate}{taux de repo} sur le Trésor
s'est établi à 5.25~\% et le taux des fonds fédéraux a dépassé sa fourchette. La Fed de New York
a prêté 53~milliards de dollars dans une opération de repo au jour le jour, le Comité a baissé ses taux
le lendemain, et à partir d'octobre elle a acheté des \omterm{def:m2:money-markets:bill}{bons du Trésor} pour reconstituer les réserves~; en 2021
elle a créé une facilité de repo permanente.
\iqlookfor{la rareté des réserves comme cause, et les trois réponses.}
\end{solution}

\begin{solution}{iq:m2:repo-and-specials:4}
Carry~: le coupon couru sur la période, diminué des intérêts du repo sur le
prix plein coupon (auquel s'ajoute une légère attraction vers le pair). Seuil de rentabilité~: la hausse de rendement qui l'annule
soit le carry divisé par le DV01~; pour un dix ans à 4.20~\% financé à 3.90~\%, environ un
quart de point de base par mois.
\iqlookfor{la formule, et la conversion en points de base via le DV01.}
\end{solution}

\begin{solution}{iq:m2:repo-and-specials:5}
Compensation multilatérale~: les transactions compensées face à la contrepartie centrale se neutralisent entre elles,
de sorte qu'un livre apparié de repo et de \omterm{def:m2:repo-and-specials:rate}{reverse repo} chez un teneur de marché se réduit à son solde net,
ce qui allège le bilan et les fonds propres~; la contrepartie centrale supprime aussi le risque
de crédit bilatéral. Le client accède au bilan des teneurs de marché à de meilleures conditions,
et l'obligation de compensation imposera de toute façon la compensation de la plupart des repo.
\iqlookfor{la compensation du bilan comme motif, et non seulement le risque de crédit.}
\end{solution}

\begin{solution}{iq:m2:repo-and-specials:6}
Oui, pour un special. Une vente à découvert qui doit livrer compare le prêt de liquidités au
\omterm{def:m2:repo-and-specials:special}{taux de repo special} $s$ pour emprunter l'obligation, qui rapporte $s$, au défaut de livraison,
qui ne rapporte rien et donne lieu à la \omterm{def:m2:repo-and-specials:fails}{pénalité de suspens} $(3\% - R)^+$. Elle emprunte tant que $s > -(3\% -
R)^+$, de sorte que les \omterm{def:m2:repo-and-specials:special}{taux de repo special} peuvent descendre jusqu'à $R - 3\%$ lorsque le taux de référence est
inférieur à 3~\%, et jusqu'à zéro dans le cas contraire. Les taux du collatéral général eux-mêmes ne deviennent
négatifs que si les taux directeurs le sont.
\iqlookfor{l'arbitrage contre le défaut de livraison, et le rôle de la \omterm{def:m2:repo-and-specials:fails}{pénalité de suspens}.}
\end{solution}
